\usepackage{algorithm}%,algpseudocode}
\usepackage{mdframed}
\usepackage{longtable}

%\cellcolor{red} % removed: cannot use \cellcolor outside a table in the preamble

% \cr is a TeX primitive used in tables - DO NOT redefine it!
% Use \mm{} (defined in packages-and-commands.tex) or \mathbf{} instead.
%\newcommand{\crbold}[1]{\mathbf{#1}} % available if needed

\newcommand{\rc}{\cellcolor{red}} % change cell color red in table
\newcommand{\hi}{\cellcolor{red}} % change cell color red in table

% Theorem environments are defined in LyryxLinAlgStyle.sty with tcolorbox styling
% Do not redefine them here to avoid "already defined" errors
%\newtheorem{example}{Example}%[section]
\newtheorem{ICE}{ICE}%[section]
%\newtheorem{problem}{Problem}%[section]  % defined in LyryxLinAlgStyle.sty
%\newtheorem{solution}{Solution}%[section]  % solutionbox defined in LyryxLinAlgStyle.sty

% Removed conflicting \input@path definitions - these are set in packages-and-commands.tex
%\def\input@path{{/figures/}}
%\def\input@path{{/path/to/folder/}{/path/to/other/folder/}}

\newcounter{algsubstate}
\renewcommand{\thealgsubstate}{\alph{algsubstate}}
\newenvironment{algsubstates}
  {\setcounter{algsubstate}{0}%
   \renewcommand{\State}{%
     \stepcounter{algsubstate}%
     \Statex {\footnotesize\thealgsubstate:}\space}}
  {}

\numberwithin{equation}{section}

\graphicspath{{figures/}{optimization/figures/figures-static/}}
\usepackage{wasysym}
\newsavebox\CBox
\newcommand\hcancel[2][0.5pt]{%
  \ifmmode\sbox\CBox{$#2$}\else\sbox\CBox{#2}\fi%
  \makebox[0pt][l]{\usebox\CBox}%
  \rule[0.5\ht\CBox-#1/2]{\wd\CBox}{#1}}

\usepackage[cc]{titlepic}
\usepackage{xcolor}
\colorlet{notgreen}{blue!50!yellow}
\newcommand{\tblue}[1]{{\color{blue} #1}}
\newcommand{\tred}[1]{{\color{red} #1}}
\newcommand{\tgreen}[1]{{\color{notgreen} #1}}

\newcommand{\code}[1]{\texttt{ #1}}
\newcommand{\complexity}[1]{{\color{red} \textit{#1}}}
\newcommand{\nphard}{\complexity{NP-Hard}}
\newcommand{\npcomplete}{\complexity{NP-Complete}}
\newcommand{\polynomial}{{\color{ansi-green} \textit{Polynomial time (P)}}}
\newcommand{\true}{\tgreen{\texttt{ true }}}
\newcommand{\false}{\tred{\texttt{ false }}}

\usepackage{scalerel}
\newcommand{\julia}{\raisebox{-.2\height}{\altincludegraphics[scale = 0.07]{julia-logo}{Julia programming language logo: three colored dots (red, green, purple) arranged in a triangle.}}  }
\newcommand{\jupyter}{\raisebox{-.2\height}{\altincludegraphics[scale = 0.09]{jupyter-logo}{Jupyter project logo: stylized orange planet circled by three gray orbital rings.}}  }
\newcommand{\jump}{\raisebox{-.2\height}{\altincludegraphics[scale = 0.08]{jump-logo}{JuMP modeling language logo: the letters 'JuMP' rendered in a bold serif typeface.}}  }
\newcommand{\gurobi}{\raisebox{-.2\height}{\altincludegraphics[scale = 0.06]{gurobi-logo}{Gurobi Optimization logo: stylized red sun-rays emblem next to the Gurobi wordmark.}}Gurobi  }
\newcommand{\coin}{\raisebox{-.2\height}{\altincludegraphics[scale = 0.2]{coin-or-logo}{COIN-OR Foundation logo: blue circular coin emblem alongside the 'COIN-OR' wordmark.}}  }
\newcommand{\ipopt}{\raisebox{-.2\height}{\texttt{Ipopt}}}

\usepackage{subcaption}

\def \PP{ {\mathcal{P}}}
\def \NN{ {\mathcal{N}}}
\def \II{ {\mathcal{I}}}
\def \ZZ{ {\mathcal{Z}}}
\def \SS{ {\mathcal{S}}}
\def \FF{ {\mathcal{F}}}
\def \CC{ {\mathcal{C}}}
\def \nn{ {\mathbb{N}}}
\def \R{ {\mathbb{R}}}
\def \Z{ {\mathbb{Z}}}
\def \Q{ {\mathbb{Q}}}
\def \cc{ {\mathbb{C}}}

\DeclareMathOperator{\dom}{dom}
\DeclareMathOperator{\cl}{cl}
\DeclareMathOperator{\intr}{intr}

\def \rank{\textup{rank}}
\def \size{\textup{size}}
%\def \dist{\textup{dist}} % defined in LyryxLinAlgStyle.sty
\def \sign{\textup{sign}}
\def \deg{\textup{deg}}
\def \conv{\textup{conv}}
\def \cone{ {\textup{cone}}}
\def \supp{\textup{supp}}
\newcommand{\interior}{ {\textup{int}}} % renamed from \int to avoid overriding the integral sign
\def \rc{\textup{rec.cone}}
\def \ri{\textup{rel.int}}
\def \rb{ {\textup{rel.bd}}}
\def \bd{ {\textup{bd}}}
\def \ls{\textup{lin.space}}
\def \tq{{\,:\,}}
\def \aff{{\textup{aff}}}
\def \one{\mathbb{1}}
\def \st{ \text{ s.t. }}

\def \BQP{\mathrm{BQP}}
\def \xLP{x^{\mathrm{LP}}}
\providecommand{\argmin}{}% ensure \argmin exists before \renewcommand
\renewcommand    \argmin         {\operatorname*{arg\,min}}
\renewcommand    \argmax         {\operatorname*{arg\,max}}

% Removed duplicate \newcounter{example} - already created by \newtheorem{example} above
%\newcounter{example}
\newcounter{general}

  % Test Environment
 \newcounter{exo}[chapter] % Reset 'exo' counter at every new section
\renewcommand{\theexo}{\thechapter.\arabic{exo}} % Format as SectionNumber.ExerciseNumber

\makeatletter
\newenvironment{exo}[1]%
{\refstepcounter{exo}%
 \protected@edef\@currentlabelname{Exercise \theexo: #1}% addition here
 \vspace{0.5cm}\noindent
 {\large\bfseries{Exercise this one \theexo~: #1} \par}
 {\par\vspace{0.5cm}}}
\makeatother

\newcounter{codeCell}
\makeatletter
\newenvironment{codeCell}%
{\refstepcounter{codeCell}%
\protected@edef\@currentlabelname{Code}% addition here
{}
%{\large\bfseries{Exercise \theexo~: #1} \par}
{}}
\makeatother

% General environment
\newenvironment{general}[2]
  {\par\medskip
   \refstepcounter{exo}%
  \protected@edef\@currentlabelname{#1}
   \begin{framed}
   \begingroup\color{black}%
   \textbf{#1: }\ignorespaces
   \par #2
   \par}
 {\endgroup\end{framed}
  \medskip}


% Example Environment
\newenvironment{optexample}
{\par\medskip
 \refstepcounter{exo}% Increment shared counter
 \protected@edef\@currentlabelname{Example \theexo}% Update label for referencing
 \begin{framed}
 \begingroup\color{blue}%
 \textbf{Example \theexo: }\ignorespaces}
{\endgroup\end{framed}
 \medskip}

% Example with code Environment
\newenvironment{examplewithcode}[2]
{\par\medskip
 \refstepcounter{exo}% Increment the shared counter and set reference anchor
 \protected@edef\@currentlabelname{Example \theexo}% Update label for referencing
 \begingroup
 \begin{center}
 \rule{\textwidth}{0.4pt}\\
 \vspace{-0.2cm}
 \rule{\textwidth}{0.4pt}
 \end{center}
 \color{black}%
 \textbf{Example \theexo: #1} % Include example number
 \hfill \protect{#2}% Code link
 \ignorespaces
 \par}
{\endgroup
 \begin{center}
 \rule{\textwidth}{1.8pt}
 \end{center}
 \medskip}

 % Define a tcolorbox-based exercise environment reusing the exo counter
% Define tcolorbox-based exercise environment


% Removed duplicate examplewithcode environment definition (was causing LaTeX error)

% Example with all code Environment
\newenvironment{examplewithallcode}[4]
{\par\medskip
 \refstepcounter{exo}% Increment shared counter
 \protected@edef\@currentlabelname{Example \theexo}% Update label for referencing
 \begin{framed}
 \begingroup\color{black}%
 \textbf{Example \theexo: #1} % Include example number
 \hfill \protect\href{#2}{\ Excel \ }\protect\href{#3}{\ PuLP \ }\protect\href{#4}{\ Gurobipy \ }% Code links
 \ignorespaces
 \par}
{\endgroup
 \end{framed}
 \medskip}

% Example without code Environment
\newenvironment{examplewithoutcode}[2]
{\par\medskip
 \refstepcounter{exo}% Increment shared counter
 \protected@edef\@currentlabelname{Example \theexo}% Update label for referencing
 \begin{framed}
 \begingroup\color{black}%
 \textbf{Example \theexo: #1} % Include example number
 \ignorespaces
 \par}
{\endgroup
 \end{framed}
 \medskip}




\newcommand{\floor}[1]{\lfloor #1 \rfloor}
\newcommand{\ceil}[1]{\lceil #1 \rceil}

\newcommand{\mathup}[1]{#1}
\newcommand{\NP}{\text{NP}}
\newcommand{\coNP}{\text{coNP}}

\newcommand{\lift}{\mathrm{lift}}
\newcommand{\ub}{\mathrm{ubbest}}
\newcommand{\F}{{\mathcal F}}

